\documentclass[11pt]{article}
\usepackage{amsmath}
\usepackage{amssymb}
\usepackage{amsfonts}
\usepackage{graphicx}
\usepackage{booktabs}
\usepackage{siunitx}
\usepackage{geometry}
\geometry{margin=1in}

\title{Paper 2: Test Case – Muon Mass from Vibrational Modes in Tetrahedral Cage \\
Standard Model Emergence in the 600-Cell Lattice Series}

\author{Thomas Lee Abshier, ND \\
Grok (xAI) \\
Hyperphysics Institute \\
https://hyperphysics.com \\
drthomas007@protonmail.com}

\date{February 1, 2026 Version 9}

\begin{document}

\maketitle

\begin{abstract}
Following repeated scaling issues in previous mass-generation attempts, this paper tests the vibrational modes approach on a single particle—the muon—as a rigorous first-principles calculation. The muon is modeled as a tetrahedral cage (4 CPs around a central core) within the 600-cell lattice. Rest mass emerges as the zero-point energy of SSV-driven harmonic oscillations. All parameters are derived from Planck scales and lattice geometry without experimental input. We provide complete step-by-step calculations, predict the muon mass, and compare to the observed value (105.66 MeV). The result guides necessary refinements and assesses the viability of the approach.
\end{abstract}

\section{Introduction}

Previous attempts at mass generation (direct SSV binding, symmetry breaking with VEV) produced unphysical scales (errors of 10^{14}–10^{26}). This paper tests the vibrational modes method on the muon (tetrahedral cage) as a controlled, transparent calculation from first principles. Success is defined as order-of-magnitude agreement (10–1000 MeV range); larger errors will pinpoint required refinements.

\section{Derivation of Restoring Force and Potential}

The restoring force for cage vibrations arises from SSV gradients between the central CP and shell CPs. For small displacements δr from equilibrium:

\begin{equation}
F = -k \delta r, \quad V(\delta r) = \frac{1}{2} k (\delta r)^2
\end{equation}

Spring constant:

\begin{equation}
k = \frac{\text{SSV}_0}{d_0^2}
\end{equation}

where SSV_0 = E_P / N_lattice, E_P = \sqrt{\hbar c^5 / G} \approx 1.22 \times 10^{28} \text{ MeV}, N_lattice = 120 \to \text{SSV}_0 \approx 1.02 \times 10^{26} \text{ MeV}, d_0 = l_P \approx 1.616 \times 10^{-35} \text{ m}.

\section{Reduced Mass and Frequency}

Reduced mass for cage of occupancy N_k = 4 (tetrahedral):

\begin{equation}
\mu_k = N_k \times m_P, \quad m_P = \sqrt{\frac{\hbar c}{G}} \approx 2.176 \times 10^{-8} \text{ kg}
\end{equation}

Angular frequency:

\begin{equation}
\omega_k = \sqrt{\frac{k}{\mu_k}} = \sqrt{\frac{\text{SSV}_0 / d_0^2}{4 \times m_P}}
\end{equation}

\section{Rest Mass from Zero-Point Energy}

Ground-state energy (QM harmonic oscillator):

\begin{equation}
E_0 = \frac{1}{2} \hbar \omega_k
\end{equation}

Rest mass:

\begin{equation}
m c^2 = E_0 = \frac{1}{2} \hbar \omega_k
\end{equation}

\section{Complete Calculation for Muon (Tetrahedral Cage)}

1. SSV_0 = 1.02 \times 10^{26} \text{ MeV}
2. d_0 = 1.616 \times 10^{-35} \text{ m}
3. k = SSV_0 / d_0^2 \approx 3.90 \times 10^{96} \text{ N/m} (in natural units)
4. \mu_k = 4 \times 2.176 \times 10^{-8} \text{ kg} \approx 8.704 \times 10^{-8} \text{ kg}
5. \omega_k = \sqrt{k / \mu_k} \approx 2.12 \times 10^{52} \text{ rad/s}
6. f_k = \omega_k / 2\pi \approx 3.37 \times 10^{51} \text{ Hz}
7. E_0 = \frac{1}{2} h f_k \approx \frac{1}{2} \times 6.626 \times 10^{-34} \times 3.37 \times 10^{51} \approx 1.12 \times 10^{18} \text{ J} \approx 7.0 \times 10^{27} \text{ MeV}

Predicted muon mass: 7.0 \times 10^{27} MeV (observed 105.66 MeV, error \sim 6.6 \times 10^{25}\%).

\section{Discussion}

The prediction is off by 26 orders of magnitude, indicating a fundamental scaling issue in the current vibrational formulation. The restoring force k is too large (Planck-scale stiffness), and the reduced mass μ_k is too small. Refinements needed:
- Stronger suppression of k (e.g., k = SSV_0 / (d_0^2 \times N_k^2)).
- Alternative reduced mass (μ_k = m_P / N_k for collective motion).
- Chiral tilt adjustment to lower effective k for higher generations.

\section{Conclusion}

This test case confirms the need for further refinement in the vibrational approach. The framework shows promise but requires additional suppression mechanisms to converge to observed scales.

\end{document}
